% == == == == == == == == == == == == == == == == == == == == == == == == == ==
    == == == == == == == == == == == ==
    = % DUAL - CHANNEL ARBITRATION % == == == == == == == == == == == == == ==
      == == == == == == == == == == == == == == == == == == == == == == == ==
    =

\newpage
\section{Dual - Channel Arbitration on the RX72N
}
\label{sec : dual - channel - arbitration}

This section documents how the RX72N firmware handles the case where data
    arrives on both USB CDC and SPI simultaneously,
    what ordering guarantees exist,
    and why the architecture prevents conflicts
            .

\subsection{The Question}

        The RX72N listens on both USB CDC and SPI at all times
            .What happens if frames arrive on both channels at the same time,
    potentially carrying different data ?

\subsection{Gateway Guarantee:
Single Active Transport
}

Under normal operation, this scenario \textbf{cannot occur} because the gateway's \texttt{TransportManager} maintains exactly one active transport:

\begin{lstlisting}[language = Go,
                        caption = {Gateway sends on ONE transport only(
                            manager.go)}] func(tm *TransportManager)
                            Send(ctx context.Context, data[] byte, ...) error {
  tm.mu.RLock() active : = tm.activeTransport // Only ONE transport
                               activeName
      : = tm.activeTransportName tm.mu.RUnlock()

              if active ==
          nil{return errors.New("no active transport available")}

          err : = active.Send(ctx, data, p...) // Send on THAT one only
                                               // ...
}
\end{lstlisting}

    The gateway \textbf{never} sends the same command on both USB and
        SPI simultaneously.It selects
        one(USB preferred, priority 10) and
    sends exclusively on that.

\subsection{When Overlap Can Occur}

        Despite the single
        - active - transport guarantee,
    brief overlap is possible during failover transitions :

\begin{enumerate}[leftmargin = 1.5em, topsep = 4pt, itemsep = 2pt]
    \item \textbf{In - flight frame during failover :
   } Frame $N$ was transmitted on USB and is still being processed by the USB
    hardware when the gateway detects a failure and switches to SPI.Frame $N
        + 1 $ is sent on SPI.The RX72N receives both frames in the same 10\,
    ms poll cycle.
    \item \textbf{Telemetry crossing :} The RX72N sends telemetry on
    USB(because the last command came from USB),
    while the gateway has already switched to SPI and sends a new command on SPI
        .The firmware receives the SPI command while USB telemetry is outgoing.
\end{enumerate}

\subsection{Sequential Polling : USB First, Then SPI}

The \texttt{rx\_comm\_manager\_poll()} function processes channels \textbf{sequentially},
    not in parallel.There is no race condition because the entire function runs
            in a single
            thread(\texttt{comm\_task}, ThreadX Priority 5, 100\, Hz)
    :

\begin{lstlisting}[language = C, caption = {Sequential channel polling(
                                      rx\_comm\_manager\_poll)}] rx_err_t
          rx_comm_manager_poll(rx_comm_manager_t * mgr) {
  bool received = false;

  /* Step 1: Poll USB channel */
  rx_err_t usb_err = internal_poll_usb(mgr);
  if (usb_err == k_rx_ok) {
    received = true; // USB frame processed
 }

  /* Step 2: Poll SPI channel */
  rx_err_t spi_err = internal_poll_spi(mgr);
  if (spi_err == k_rx_ok) {
    received = true; // SPI frame processed
 }

  /* Step 3: Process event queue */
  internal_process_event_queue(mgr);

  /* Step 4: Check heartbeat timers */
  internal_check_heartbeat(mgr);

  return received ? k_rx_ok : k_rx_err_timeout;
}
\end{lstlisting}

\textbf{Ordering guarantee :} Within any single poll cycle,
    USB is always processed before SPI.This is deterministic and repeatable.

\subsection{Frame Processing Pipeline}

        Both \texttt{
            internal\_poll\_usb()} and \texttt{internal\_poll\_spi()} follow the
            same pipeline :

\vspace {
  0.3cm
}

\noindent\makebox[\textwidth] {
  %
\begin{tikzpicture}[font =\small, node distance=0.3cm,
                     box/.style={draw, thick, rectangle, minimum width = 4cm,
                                    minimum height = 0.8cm, align = center,
                                    rounded corners = 2pt},
                     arrow/.style={->, thick, >= stealth}] %
      USB path
    \node[box, fill=blue!15](usb_recv)
          at(0, 0){USB : rx\_usb\_comm\_receive()};
  \node[box, fill=yellow!15, below=of usb_recv](usb_handle){
      internal\_handle\_frame(USB)};
  \node[box, fill=green!15, below=of usb_handle](usb_cb){
      internal\_frame\_callback()};
  \node[box, fill=orange!15, below=of usb_cb](usb_cmd){
      shared\_data\_set\_motor\_command()};

  % SPI path
    \node[box, fill=red!15](spi_recv)
          at(7, 0){SPI : rx\_spi\_link\_receive()};
  \node[box, fill=yellow!15, below=of spi_recv](spi_handle){
      internal\_handle\_frame(SPI)};
  \node[box, fill=green!15, below=of spi_handle](spi_cb){
      internal\_frame\_callback()};
  \node[box, fill=orange!15, below=of spi_cb](spi_cmd){
      shared\_data\_set\_motor\_command()};

  % Arrows
    \draw[arrow](usb_recv)--(usb_handle);
  \draw[arrow](usb_handle)--(usb_cb);
  \draw[arrow](usb_cb)--(usb_cmd);
  \draw[arrow](spi_recv)--(spi_handle);
  \draw[arrow](spi_handle)--(spi_cb);
  \draw[arrow](spi_cb)--(spi_cmd);

  % Time arrow
    \draw[->, very thick, gray](-2.5, 0.5)--(
        -2.5, -4)node[below, font =\footnotesize]{Time};

  % Labels
    \node[above, font =\footnotesize\bfseries] at(usb_recv.north){
        Processed First};
  \node[above, font =\footnotesize\bfseries] at(spi_recv.north){
      Processed Second};

  % Overwrite indicator
    \draw[->, thick, red, dashed](
        spi_cmd.west)-- node[below, font =\footnotesize\color{red}]{
        Overwrites}(usb_cmd.east);
  \end{tikzpicture} %
}

\vspace { 0.3cm}

\subsection{Last Writer Wins}

Both frames call \texttt{shared\_data\_set\_motor\_command(\& cmd)},
    which is mutex -
        protected.Since the two calls execute sequentially(not concurrently),
    the second call \textbf{overwrites} the first :

\begin{enumerate}[leftmargin = 1.5em, topsep = 4pt, itemsep = 2pt]
    \item USB frame arrives $\rightarrow$ \texttt{
        shared\_data\_set\_motor\_command()} writes motor
    velocities from USB frame.
    \item SPI frame arrives $\rightarrow$ \texttt {
  shared\_data\_set\_motor\_command()
}
\textbf{overwrites} with motor velocities from SPI frame.
    \item Motor task wakes up $\rightarrow$ reads the \textbf{SPI frame's data} (the last one written).
\end{enumerate}

\textbf{The SPI command wins because it is polled second.} The USB command is effectively superseded.

\subsection{Why This Is Correct During Failover}

During a failover, the overlap works \textbf{correctly} because the gateway uses shared \texttt{SessionState} sequence numbers:

\begin{enumerate}[leftmargin=1.5em,topsep=4pt,itemsep=2pt]
    \item The gateway increments the sequence number for each \texttt{Send()} call (shared counter).
    \item Frame $N$ goes out on USB (the old transport).
    \item Frame $N+1$ goes out on SPI (the new transport).
    \item The RX72N processes USB ($N$) first, then SPI ($N+1$) second.
    \item The motor task gets frame $N+1$, which is the \textbf{newer} command.
    \item This is the correct behavior --- the motor should execute the most recent command.
\end{enumerate}

\noindent\colorbox{keyinsightcolor} {
  \parbox{\dimexpr\linewidth - 2\fboxsep} {
    %
\textbf{Key Insight :
   } The ``last writer wins'' semantics combined with sequential
             polling(USB before SPI) and
        monotonically increasing sequence numbers guarantee that the motor task
            always receives the most recent command,
        even during transport failover.
 }
}

\subsection{Sequence Validation}

The shared \texttt{
    rx\_session} module on the RX72N validates every received sequence number :

\begin{itemize}[leftmargin = 1.5em, topsep = 4pt, itemsep = 2pt]
    \item \textbf{Duplicate detection :} If the same sequence number is received
    twice (e.g., frame retransmitted on both channels),
    the second copy is rejected by \texttt{rx\_session\_validate\_rx()}.
    \item \textbf{Gap tolerance :} Sequence gaps of up to 10 frames are accepted
    (see Section ~\ref{subsec : shared - session}).
    This handles frames lost during failover.
    \item \textbf{Wraparound:
}
16 - bit sequence numbers wrap from 65535 to 0. The validation logic
         handles this correctly.
\end{itemize}

\subsection{Channel Parameter(Currently Unused)}

    The \texttt{internal\_handle\_command\_frame()} function receives a \texttt{
        channel} parameter indicating whether the frame arrived on USB or
    SPI :

\begin{lstlisting}[language = C,
                    caption = {Channel parameter in command handler(
                        internal\_handle\_velocity\_request)}] static void
    internal_handle_command_frame(rx_comm_channel_t channel,
                                  const rx_frame_t *frame) {
  (void)channel; // Currently unused
                 // ... decode protobuf, update shared_data ...
}
\end{lstlisting}

The channel parameter is currently cast to \texttt{void} (unused). A future enhancement could use it to:
\begin{itemize}[leftmargin=1.5em,topsep=4pt,itemsep=2pt]
    \item Route response/telemetry back on the same channel that sent the command.
    \item Log which channel is actively receiving commands for diagnostics.
    \item Implement channel-specific command filtering (e.g., only accept emergency stop from USB).
\end{itemize}

\subsection{Communication Watchdog Interaction}

The 500\,
    ms communication watchdog(\texttt{shared\_data\_update\_last\_comm\_tick()})
        is updated on \textbf{every valid frame from either channel} :

\begin{lstlisting}[language = C, caption = {Watchdog update on any frame(
                                      internal\_update\_watchdog)}] static void
            internal_frame_callback(rx_comm_channel_t channel,
                                    const rx_frame_t *frame, void *ctx) {
  (void)ctx;
  if (frame == nullptr) {
    return;
 }

  /* Update watchdog on ANY valid frame from ANY channel */
  shared_data_update_last_comm_tick();

  switch (frame->header.type) {
  case k_frame_type_command:
    internal_handle_command_frame(channel, frame);
    break;
    // ...
 }
}
\end{lstlisting}

This means :
\begin{itemize}[leftmargin = 1.5em, topsep = 4pt, itemsep = 2pt]
    \item During failover,
    if USB stops sending but SPI starts,
    the watchdog is still satisfied.
    \item The 500\,
    ms timer only triggers emergency stop if \textbf{both} channels are completely silent
        .
    \item A single valid frame on either channel resets
            the watchdog and prevents emergency stop.
\end{itemize}

\subsection{Summary:
  Dual - Channel Behavior Matrix
}

\begin{table}[H]
\centering
\caption {RX72N Behavior for Dual-Channel Scenarios
}
\label {
tab:
  dual - channel - matrix
}
\begin{tabular} { l l l}
\toprule
\textbf{Scenario} & \textbf{What Happens} & \textbf{Result} \\
\midrule
Normal operation    & Gateway sends on ONE channel only          & Single frame per cycle \\
Failover overlap    & Both channels have data                   & USB processed first, SPI overwrites \\
                    &                                           & (newer command = correct) \\
Duplicate frame     & Same sequence on both channels             & Sequence validation rejects second \\
Different senders   & Cannot happen (single gateway)            & Architecture prevents this \\
All channels silent & No frames for 500\,ms                     & Emergency stop triggered \\
One channel fails   & Other channel still delivers frames       & Watchdog still satisfied \\
\bottomrule
\end{tabular}
\end{table}

\subsection{Cross - References}

\begin{itemize}[leftmargin = 1.5em, topsep = 4pt, itemsep = 2pt]
    \item \textbf{Section ~\ref{sec : communication - stack}} : Protocol stack,
                                                                 frame format,
                                                                 CRC validation
    \item \textbf{Section ~\ref{sec : transport - failover}}
    : Transport failover,
      health monitoring,
      shared SessionState
    \item \texttt{e2 - studio - star - rx72n -
                   firmware / libs / rx\_comm\_manager / src /
                       rx\_comm\_manager.c}
    : Dual -
      channel polling implementation
    \item \texttt{e2 - studio - star - rx72n -
                   firmware / src / tasks / comm\_task.c} : Frame callback,
      command dispatch,
      watchdog update
    \item \texttt{star - gateway / internal / manager / manager.go}
    : Single -
      active - transport guarantee
\end{itemize}
